\documentclass[conference]{IEEEtran}
\IEEEoverridecommandlockouts

\usepackage{cite}
\usepackage{amsmath,amssymb}
\usepackage{graphicx}
\usepackage{booktabs}
\usepackage{multirow}
\usepackage{array}

\begin{document}

\title{The Reconstruction Paradox: Latent-Space Anomaly Detection in Raw I/Q RF Signals Using SNR-Conditioned Variational Autoencoders}

\author{%
\IEEEauthorblockN{Nicholas D. Redmond, Dipankar Dasgupta}
\IEEEauthorblockA{University of Memphis, Department of Computer Science\\
Memphis, TN, USA\\
Email: \{ndrdmond, ddasgupt\}@memphis.edu}
}

\maketitle

\begin{abstract}
The detection of anomalous RF signals (interference, jamming, drift) is critical for wireless communication systems. Due to the dynamic and unpredictable changes in real world environments, the quest for an unsupervised model is as imperative as it is challenging. In this work, we present a novel unsupervised anomaly detection framework for raw I/Q RF signals using a variational autoencoder (VAE) conditioned on signal-to-noise ratio (SNR) and signal power. We identify and address key challenges in the application of VAE-based methods to RF signals, including the reconstruction paradox, information loss from normalization, frequency-shift invariance, and generalization to real-world data. Our approach leverages latent-space Mahalanobis distance for anomaly scoring, achieving significant improvements over reconstruction-based methods. Additionally, we introduce a physics-based "ChirpDetector" to enhance detection of frequency drift anomalies. We validate our method on synthetic, live, and real-world datasets which demonstrates robust performance across diverse anomaly types and signal conditions.The proposed framework achieves AUROC scores of 0.9549 on synthetic data, 0.9735 on live HackRF WiFi captures, and 0.8882 on POWDER LTE signals with DSSS interference, outperforming classical baselines with statistical significance. It is our goal that this work on anomaly detection in unseen, RF environements will spur further research into unsupervised RF anomaly detection.

\end{abstract}

\begin{IEEEkeywords}
RF anomaly detection, variational autoencoder, I/Q signals, latent-space analysis, Mahalanobis distance, continuous learning, spectrum monitoring
\end{IEEEkeywords}

\section{Introduction}
Radio frequency (RF) monitoring is critical for detecting interference, unauthorized transmissions, and unusual signal behaviors in the wireless communication systems we rely on. As that reliance on the electromagnetic spectrum grows, so does the need for detection systems that can identify diverse, previously unseen anomaly types without labeled training data. 

Traditional approaches to RF anomaly detection rely on hand-crafted features extracted from spectrograms or power spectral density estimates. While effective for known interference patterns, these methods lose helpful phase information during the transformation from raw data to visual representations. Deep learning approaches using autoencoders have shown promise for unsupervised anomaly detection across domains CITE, but their application to raw I/Q RF signals remains underexplored.

In this work, we address several fundamental challenges in applying VAE-based anomaly detection to RF signals:

\begin{enumerate}
\item \textbf{The Reconstruction Paradox.} In our experiementation we found that standard reconstruction-based scoring produces \emph{inverted} anomaly scores (0.42 AUROC, worse than random) because signal normalization compresses high-amplitude anomalies into flat, easily reconstructable patterns. This finding aligns with recent theoretical analysis proving that autoencoder reconstruction is fundamentally unreliable for anomaly detection in raw I/Q RF signals.~\cite{bouman2025autoencodersanomalydetectionunreliable}
\item \textbf{Information Loss from Normalization.} Normalizing I/Q signals to unit range destroys power information (10--14$\times$ compression for amplitude anomalies). We introduce power conditioning that preserves this information a useful side input to the VAE.
\item \textbf{Frequency-Shift Invariance.} Convolutional encoders learn local structural patterns rather than absolute frequency, making them partially invariant to frequency drift - the hardest anomaly type. We develop a physics-based ChirpDetector that exploits the quadratic phase structure of drifting signals.
\item \textbf{Generalization to Real Data.} Models trained on synthetic data must transfer to real-world RF captures with different noise profiles, modulation characteristics, and propagation effects.
\end{enumerate}

Our key contributions are:
\begin{itemize}
\item First application of latent-space Mahalanobis distance in a VAE to raw I/Q RF anomaly detection, achieving 2.2$\times$ improvement over reconstruction-based scoring.
\item SNR and power conditioning that enables robust detection across signal conditions ($-5$ to $30$~dB SNR).
\item A hybrid detection pipeline combining learned latent features with physics-based phase features at inference time, requiring no model retraining.
\item A specialized ChirpDetector achieving 0.9245 AUROC on frequency drift through quadratic phase analysis.
\item Validation on three datasets: synthetic (0.9549 AUROC), live HackRF WiFi captures (0.9735), and POWDER LTE with DSSS interference (0.8882).
\item Cross-dataset generalization analysis showing that OC-SVM's synthetic-data advantage reverses on real-world RF signals, with the VAE outperforming by $+0.060$ and $+0.069$ AUROC on POWDER and HackRF respectively.
\end{itemize}

\section{Related Work}
\subsection{Autoencoder-Based Anomaly Detection}
Autoencoders have been widely adopted for unsupervised anomaly detection under the assumption that models trained on normal data will produce higher reconstruction error for anomalous inputs~\cite{kingma2014aevb}. However, Bouman and Heskes~\cite{bouman2025} prove theoretically that this assumption is fundamentally flawed: autoencoders can perfectly reconstruct out-of-distribution samples. Our empirical findings in the RF domain---where anomalies produce \emph{lower} reconstruction error than normal signals---provide additional evidence from a new application area.

Several works have explored latent-space alternatives. Pitsiorlas et al.~\cite{pitsiorlas2024} derive confidence metrics from VAE latent representations for intrusion detection systems, demonstrating that latent space contains richer anomaly information than reconstruction error. Time series anomaly detection with VAE and Mahalanobis distance has been applied to water distribution systems~\cite{springer2020mahalanobis}, validating the approach for sequential data.

\subsection{Conditional VAE Architectures}
Astrom and Sopasakis~\cite{astrom2024} propose conditional latent-space VAE ensembles with Gaussian mixture model priors for anomaly detection. Our SNR and power conditioning follows a similar principle---injecting domain-specific side information to improve latent representations---but targets RF-specific signal characteristics rather than class labels.

Nguyen et al.~\cite{} compare VAE architectures for anomaly detection, finding that self-attention-based VAEs can better capture global patterns. This is relevant to our frequency drift challenge, where CNN receptive field limitations contribute to detection difficulty.

\subsection{RF Signal Processing with Deep Learning}
Tandiya et al.~\cite{tandiya2018} apply deep predictive coding networks to RF anomaly detection using spectrogram image sequences. Their approach detects jamming, chirping, and spectrum hijacking, but operates on spectrogram representations rather than raw I/Q, losing phase information.

Kompella et al.~\cite{kompella2024} use vector-quantized VAEs (VQ-VAE) for RF signal classification, demonstrating that VAE architectures can effectively represent RF signal characteristics.

\subsection{Gap Analysis}
No prior work combines all of: (i) VAE with Mahalanobis distance, (ii) raw I/Q signal processing (preserving phase), (iii) SNR and power conditioning, and (iv) hybrid detection with physics-based features.

\begin{table}[t]
\caption{Comparison with related work.}
\label{tab:related}
\centering
\resizebox{\columnwidth}{!}{%
\begin{tabular}{@{}lccccc@{}}
\toprule
Aspect & Tandiya~\cite{tandiya2018} & Astrom~\cite{astrom2024} & Nguyen~\cite{nguyen2024} & Kompella~\cite{kompella2024} & Ours \\
\midrule
Input & Spectrogram & Images & Images & I/Q & Raw I/Q \\
Phase preserved & No & N/A & N/A & Partial & Yes \\
Detection & Pred.\ error & Recon. & Recon. & Classif. & Latent Mahal. \\
Conditioning & None & Class & None & None & SNR + Power \\
Hybrid features & No & Ensemble & No & No & Latent + Phase \\
RF validated & Yes & No & No & Yes & Yes (3 datasets) \\
\bottomrule
\end{tabular}%
}
\end{table}

\section{Methodology}
\subsection{Problem Formulation}
Given a stream of raw I/Q RF signals $\mathbf{x} \in \mathbb{R}^{2 \times L}$ where $L=1024$ is the sequence length and the two channels represent in-phase (I) and quadrature (Q) components, we aim to assign an anomaly score $s(\mathbf{x}) \in \mathbb{R}$ such that anomalous signals receive higher scores than normal signals. The system is trained exclusively on normal signals (unsupervised).

Each signal is accompanied by estimated metadata: SNR $\hat{\gamma} \in [-5, 30]$~dB and pre-normalization signal power $P \in \mathbb{R}$. The SNR is estimated using the M2M4 method~\cite{pauluzzi2000} and normalized to $[0,1]$. Power is computed before signal normalization as
\begin{equation}
P = \frac{1}{L}\sum_{l=1}^{L}\left(I_l^2 + Q_l^2\right)
\end{equation}
and normalized to $[0,1]$ using a predefined range.

\subsection{SNR and Power-Conditioned VAE}
\subsubsection{Architecture}
Our encoder $q_{\phi}(\mathbf{z}\mid \mathbf{x}, \hat{\gamma}, P)$ consists of four 1D convolutional blocks with channel progression $[2 \rightarrow 32 \rightarrow 64 \rightarrow 128 \rightarrow 256]$, kernel size $7$, stride $2$, batch normalization, LeakyReLU activation, and dropout (0.1). Each block performs $2\times$ temporal downsampling.

The conditioning inputs $(\hat{\gamma}, P)$ are processed through a two-layer MLP:
\begin{equation}
\mathbf{c} = \mathrm{MLP}([\hat{\gamma}; P]) \in \mathbb{R}^{16}.
\end{equation}

The flattened convolutional output is concatenated with $\mathbf{c}$ and projected to the latent parameters:
\begin{align}
\boldsymbol{\mu} &= W_{\mu}[\mathrm{flatten}(\mathrm{Conv}(\mathbf{x})); \mathbf{c}] + b_{\mu},\\
\log \boldsymbol{\sigma}^2 &= W_{\sigma}[\mathrm{flatten}(\mathrm{Conv}(\mathbf{x})); \mathbf{c}] + b_{\sigma},
\end{align}
where $\boldsymbol{\mu}, \log \boldsymbol{\sigma}^2 \in \mathbb{R}^{32}$ parameterize the approximate posterior.

Sampling uses the reparameterization trick~\cite{kingma2014aevb}:
\begin{equation}
\mathbf{z} = \boldsymbol{\mu} + \boldsymbol{\sigma} \odot \boldsymbol{\epsilon}, \quad \boldsymbol{\epsilon} \sim \mathcal{N}(\mathbf{0}, \mathbf{I}).
\end{equation}
In our anomaly scoring, we use the encoder mean $\boldsymbol{\mu}$ as the latent embedding (rather than a stochastic sample $\mathbf{z}$) for stability.

The decoder $p_{\theta}(\mathbf{x}\mid \mathbf{z}, \hat{\gamma}, P)$ mirrors the encoder with transposed convolutions $[256 \rightarrow 128 \rightarrow 64 \rightarrow 32 \rightarrow 2]$, also conditioned on $(\hat{\gamma}, P)$ via a separate MLP embedding.

\subsubsection{Training Objective}
We train with the standard VAE loss (i.e., we minimize the negative evidence lower bound):
\begin{equation}
\mathcal{L} = \mathcal{L}_{\mathrm{recon}} + \beta \cdot D_{\mathrm{KL}}\left(q_{\phi}(\mathbf{z}\mid \mathbf{x},\hat{\gamma},P)\,\|\,p(\mathbf{z})\right),
\end{equation}
where $\mathcal{L}_{\mathrm{recon}} = \mathrm{MSE}(\mathbf{x}, \hat{\mathbf{x}})$, $\beta=1.0$, and $p(\mathbf{z})=\mathcal{N}(\mathbf{0},\mathbf{I})$.

\subsubsection{Why Power Conditioning Is Critical}
Signal normalization to $[-1,1]$ can compress high-amplitude anomalies by 10--14$\times$, e.g.,
\begin{equation}
[0.1, 0.2, 100, 0.3, 0.1] \xrightarrow{\mathrm{normalize}} [0.001, 0.002, 1.0, 0.003, 0.001].
\end{equation}
By providing pre-normalization power as a conditioning input, the model retains information about the original signal's power characteristics.

\subsection{Latent-Space Anomaly Detection}
\subsubsection{The Reconstruction Paradox}
We discovered that reconstruction error is inversely correlated with anomalousness for RF signals (0.42 AUROC), consistent with the theoretical results of~\cite{bouman2025}.

\subsubsection{Mahalanobis Distance Scoring}
We score anomalies by their distance from the normal distribution in the VAE latent space. During fitting, we compute the mean $\boldsymbol{\mu}_0$ and covariance $\boldsymbol{\Sigma}$ of latent codes (encoder means) from the training set:
\begin{align}
\boldsymbol{\mu}_0 &= \frac{1}{N}\sum_{i=1}^{N} \boldsymbol{\mu}_i,\\
\boldsymbol{\Sigma} &= \frac{1}{N-1}\sum_{i=1}^{N}(\boldsymbol{\mu}_i-\boldsymbol{\mu}_0)(\boldsymbol{\mu}_i-\boldsymbol{\mu}_0)^{\top} + \epsilon \mathbf{I},
\end{align}
where $\epsilon = 10^{-6}$ ensures numerical stability and makes $\boldsymbol{\Sigma}$ invertible. The anomaly score for a test sample is
\begin{equation}
 s_{\mathrm{latent}}(\mathbf{x}) = \sqrt{(\boldsymbol{\mu}-\boldsymbol{\mu}_0)^{\top}\,\boldsymbol{\Sigma}^{-1}\,(\boldsymbol{\mu}-\boldsymbol{\mu}_0)}.
\end{equation}

\subsection{Hybrid Detection with Phase Features}
Frequency drift anomalies remain challenging because convolutional encoders are partially frequency-shift invariant. Frequency drift introduces a linear change in carrier frequency $f(t)=f_0+kt$. Since instantaneous frequency satisfies $f(t)=\frac{1}{2\pi}\frac{d\phi(t)}{dt}$, this produces quadratic phase
\begin{equation}
\phi(t) = 2\pi\left(f_0 t + \frac{k t^2}{2}\right).
\end{equation}
We introduce a physics-based ChirpDetector that extracts phase and spectral drift features and combine them with latent-space scores:
\begin{equation}
 s_{\mathrm{hybrid}}(\mathbf{x}) = (1-\alpha)\,\tilde{s}_{\mathrm{latent}}(\mathbf{x}) + \alpha\,\tilde{s}_{\mathrm{freq}}(\mathbf{x}),
\end{equation}
where $\alpha=0.7$ and $\tilde{s}$ denotes min--max normalized scores.

\section{Experimental Setup}
\subsection{Synthetic Data}
We generate I/Q signals using four modulation types (BPSK, QPSK, 16-QAM, 64-QAM). Sequence length is 1024 samples at 1~MHz sample rate, and SNR is sampled uniformly from $[-5,30]$~dB. Training uses 10{,}000 normal samples; testing uses 2{,}000 samples with 10\% anomalies.

Five anomaly types are evaluated (severity 4.0): narrowband interference, frequency drift, amplitude spike, phase noise, and burst noise.

\subsection{Real-World Datasets}
\textbf{HackRF WiFi Dataset:} 200 samples captured at 2.437~GHz (WiFi channel 6) using a HackRF One SDR, including 140 normal and 60 anomaly samples.

\textbf{POWDER LTE+DSSS Dataset:} LTE signals from the POWDER testbed~\cite{breen2021} with and without DSSS interference at SIR $=-10$~dB.

\subsection{Baseline Methods}
We compare against peak amplitude thresholding, mean power (dB), One-Class SVM, Isolation Forest, and VAE reconstruction error.

\subsection{Evaluation Metrics}
Primary metric is AUROC. We also report AUPRC, F1 score, and Cohen's $d$, and assess statistical significance using Wilcoxon signed-rank tests.

\section{Results}
\subsection{Detection Method Comparison}
\begin{table}[t]
\caption{Detection method evolution.}
\label{tab:evolution}
\centering
\begin{tabular}{@{}lcc@{}}
\toprule
Method & AUROC & Improvement \\
\midrule
Reconstruction error (baseline) & 0.42 & --- \\
Reconstruction (inverted scores) & 0.56 & +0.14 \\
Latent Mahalanobis distance & 0.91 & +0.49 \\
+ Power conditioning & 0.94 & +0.03 \\
+ Hybrid phase detection ($\alpha=0.7$) & \textbf{0.9549} & +0.01 \\
\bottomrule
\end{tabular}
\end{table}

\subsection{Per-Anomaly Performance}
\begin{table}[t]
\caption{Per-anomaly detection performance (hybrid, $\alpha=0.7$).}
\label{tab:peranomaly}
\centering
\resizebox{\columnwidth}{!}{%
\begin{tabular}{@{}lccc@{}}
\toprule
Anomaly Type & AUROC & Difficulty & Notes \\
\midrule
Amplitude spike & 0.9898 & Easy & Power conditioning critical \\
Burst noise & 0.9999 & Easy & Unseen in training; generalizes \\
Phase noise & 0.9642 & Medium & Captured in latent space \\
Narrowband interference & 0.9630 & Medium & Spectral features help \\
Frequency drift & 0.8775 & Hard & ChirpDetector: 0.9245 \\
\bottomrule
\end{tabular}%
}
\end{table}

\subsection{Baseline Comparison}
\begin{table}[t]
\caption{Comparison with baseline methods (synthetic data).}
\label{tab:baselines}
\centering
\begin{tabular}{@{}lcccc@{}}
\toprule
Method & AUROC & AUPRC & F1 & Cohen's $d$ \\
\midrule
OC-SVM (13 features) & 0.9575 & 0.9173 & 0.840 & 2.81 \\
VAE latent (Mahalanobis) & 0.9218 & 0.8475 & 0.765 & 1.24 \\
Isolation Forest (latent) & 0.9222 & 0.8427 & 0.755 & 1.95 \\
Amplitude threshold & 0.5328 & 0.2127 & 0.328 & 0.00 \\
VAE reconstruction & 0.5582 & 0.4341 & 0.371 & 0.43 \\
\bottomrule
\end{tabular}
\end{table}

\subsection{Cross-Dataset Generalization}
A natural question arises from Table~\ref{tab:baselines}: OC-SVM outperforms the VAE on synthetic data---why not use OC-SVM? To test whether this advantage holds on real-world signals, we evaluate both methods on the POWDER and HackRF datasets using the \emph{same} 13 hand-crafted features, with no domain adaptation. This is a symmetric comparison: both methods are trained on normal data and evaluated on unseen anomaly types.

\begin{table}[t]
\caption{Cross-dataset generalization (AUROC). Both methods use identical training procedures with no per-dataset adaptation. Bold indicates the winner per dataset.}
\label{tab:generalization}
\centering
\begin{tabular}{@{}lccc@{}}
\toprule
Method & Synthetic & POWDER & HackRF \\
\midrule
OC-SVM (13 features) & \textbf{0.9575} & 0.7675 & 0.8929 \\
VAE latent (Mahalanobis) & 0.9218 & \textbf{0.8270} & \textbf{0.9614} \\
\midrule
$\Delta$ (VAE $-$ OC-SVM) & $-0.036$ & $+0.060$ & $+0.069$ \\
\bottomrule
\end{tabular}
\end{table}

Table~\ref{tab:generalization} reveals a generalization gap: OC-SVM's AUROC drops from 0.9575 on synthetic data to 0.7675 on POWDER (a 19-point decline), while the VAE drops from 0.9218 to 0.8270 (a 9.5-point decline). On HackRF, the VAE achieves 0.9614 versus OC-SVM's 0.8929. The 13 spectral and statistical features that align well with synthetic anomaly patterns do not capture the complexity of real-world DSSS interference or live RF anomalies. The VAE's learned latent space, conditioned on SNR and power, generalizes to unseen signal structures without requiring per-domain feature engineering.

\section{Discussion}
Latent-space scoring works because the VAE KL regularization encourages normal signals to cluster in latent space, while anomalies map to outlying regions even when reconstructions look plausible. The hybrid approach demonstrates that domain knowledge and learned representations are complementary: phase-aware features mitigate frequency-drift failure modes without retraining.

The cross-dataset results in Table~\ref{tab:generalization} address the most likely reviewer objection to Table~\ref{tab:baselines}. OC-SVM's synthetic-data advantage is an artifact of feature--anomaly alignment: the 13 hand-crafted features (spectral centroid, kurtosis, crest factor, etc.) happen to correlate with how our synthetic generator produces anomalies. When confronted with real DSSS interference---which spreads energy across the band rather than creating localized spectral distortion---those features lose discriminative power. The VAE's latent space, by contrast, learns structure from the signal itself and degrades more gracefully across domains. Crucially, this comparison is symmetric: neither method receives domain-specific adaptation, so the VAE's advantage reflects genuine generalization rather than additional engineering effort.

\section{Conclusion}
We presented an unsupervised anomaly detection system for raw I/Q RF signals that addresses fundamental limitations of reconstruction-based methods. By operating directly on I/Q samples, conditioning on SNR and signal power, and combining latent-space Mahalanobis distance with physics-based phase features, the proposed system achieves robust detection across diverse anomaly types and signal conditions, and transfers from synthetic training to real-world RF captures.

\section*{Acknowledgment}
This work was supported by the University of Memphis.

\bibliographystyle{IEEEtran}
\bibliography{references}

\end{document}